\section{Methodology}

\subsection{Profile Architecture}

Each reviewer profile is a markdown file with YAML frontmatter containing:

\begin{itemize}
\item \textbf{Metadata}: Name, affiliation, category, keywords
\item \textbf{Research Background}: 2-3 sentence overview of expertise
\item \textbf{Key Publications}: 3-5 representative papers with impact
\item \textbf{Evaluation Lens}: Characteristic questions the reviewer asks
\item \textbf{Review Criteria}: Structured checklist aligned with research priorities
\item \textbf{Voice \& Tone}: Style guide for generating reviews in this persona
\end{itemize}

Profiles average 2KB each, stored in \texttt{context/panel/reviewers/profiles/\{name\}.md}.

\subsection{Resolution Chain}

The profile loader implements a three-tier fallback:

\begin{enumerate}
\item \textbf{Exact match}: \texttt{profiles/percy-liang.md}
\item \textbf{Slug match}: Convert "Percy Liang" $\rightarrow$ \texttt{percy-liang.md}
\item \textbf{Database fallback}: Extract from \texttt{REVIEWER-DATABASE.md} (current behavior)
\end{enumerate}

Session-level caching prevents repeated file system access within a single review cycle.

\subsection{Experimental Design}

We conduct a controlled A/B comparison with 5 test papers across varied venues (CHI, NeurIPS, ICML, ACL, PLDI). Each paper undergoes:

\begin{itemize}
\item \textbf{Condition A (Baseline)}: Current system with database parsing
\item \textbf{Condition B (Profiles)}: Profile loader with caching
\end{itemize}

Instrumentation captures total tokens, reviewer context tokens, and per-round breakdown. Reviews are generated using identical prompts aside from persona delivery method.

\subsection{Metrics}

\begin{itemize}
\item \textbf{Token Efficiency}: Reduction percentage in reviewer context tokens
\item \textbf{Quality Preservation}: Cosine similarity of P1 items between conditions
\item \textbf{Consistency}: Cross-round correlation of reviewer positions (Round 1 $\leftrightarrow$ Round 2)
\item \textbf{Diversity}: Inter-reviewer perspective variation (manual annotation)
\end{itemize}
